\documentclass[11pt,a4paper,DIV9,german]{scrartcl}
%%%%%%%%%%%%%%%%%%%%%%%%%%%%%%%%%%%%%%%%%%%%%%%%%%%%%%%%%%%%%%%%%%%%%%%%%%
% Anpassungen f"ur das jeweilige "Ubungsblatt
\def\blnr{1}              % Blattnummer
\def\sdate{01.03.2019,}     % Abgabedatum
\def\exnr{1}              % Nummer der ersten Aufgabe des aktuellen Blattes
%%%%%%%%%%%%%%%%%%%%%%%%%%%%%%%%%%%%%%%%%%%%%%%%%%%%%%%%%%%%%%%%%%%%%%%%%%
\usepackage{babel,theorem}
\usepackage{exscale}
\usepackage{amsmath,amsfonts,amssymb}
\usepackage{a4wide}
\usepackage{epic,eepic}
\usepackage{graphicx} % einbinden von Grafiken mit \includegraphics
\usepackage{url}

\pagestyle{empty}
\parindent0em
\parskip1ex
\addtolength{\topmargin}{0cm}
\addtolength{\oddsidemargin}{0cm}
\addtolength{\textheight}{2cm}
\addtolength{\textwidth}{0cm}

% neuer Z"ahler f"ur die Aufgaben
\newcounter{aufgnr}
\setcounter{aufgnr}{\exnr}
\addtocounter{aufgnr}{-1}

% in enumerate-Umgebungen erst Kleinbuchstaben, dann kleine r"omische Zahlen
\renewcommand{\labelenumi}{(\alph{enumi})\,}
\renewcommand{\labelenumii}{(\roman{enumii})\,}

% Stil f"ur Definitionen, S"atze usw.
\theoremstyle{break}   % Titelzeile abgesetzt
\newtheorem{Satz}{Satz}
\newtheorem{Lem}{Lemma}
\newtheorem{Def}{Definition}
\theorembodyfont{\upshape}  % normaler Font f"ur Beispiel, Aufgabe
\newtheorem{Bsp}{Beispiel}
\newtheorem{aufg}[aufgnr]{Aufgabe}
\newenvironment{Aufgabe}[1][z*+]%
{\begin{aufg}}
  {\end{aufg}\vspace{1.5ex}}
\newtheorem{praktaufg}[aufgnr]{Aufgabe^*}
\newenvironment{PraktAufgabe}[1][]%  der optionale Param. muss in [[ .. ]]
{\begin{praktaufg}#1}
  {\end{praktaufg}\vspace{1.5ex}}

\newcommand{\zweifaelles}[3]{ \left\{ \begin{array}{ccc} #1 & \hspace{.8cm}  \mbox{ falls } #2 \, , &     \\[1ex] 
                               #3 & \mbox{ sonst } &     \end{array} \right.     }
% Mathematische K"urzel
\def\R{\mathbb{R}}
\def\S{\mathbb{S}}
\def\C{\mathbb{C}}
\def\Q{\mathbb{Q}}
\def\N{\mathbb{N}}
\def\Z{\mathbb{Z}}
\def\P{\mathbb{P}}
\def\K{\mathbb{K}}
\def\O{{\mathcal O} }
\def\G{{\mathcal G} }
\def\T{{\mathcal T} }
\def\U{{\mathcal U} }

\newcommand{\arrl}{\begin{array}{l}}
\newcommand{\are}{\end{array}}
\newcommand{\arrc}{\begin{array}{c}}
\newcommand{\mza}{\begin{eqnarray}}
\newcommand{\mze}{\end{eqnarray}}
\newcommand{\dn}{\frac{\partial}{\partial n}}
\newcommand{\dny}{\frac{\partial}{\partial n(y)}}
\newcommand{\dd}[2]{\ensuremath{\frac{\partial^{#1}}{\partial #2}}} 
\newcommand{\ov}{\overline}
\newcommand{\ma}{\mid\!}
\newcommand{\me}{\!\mid}
\newcommand{\impl}{\Rightarrow}
\newcommand{\sca}[3]{\ensuremath{(#1,#2)_{#3}}}
\newcommand{\nor}[2]{\ensuremath{\parallel\! #1 \!\parallel_{#2}}}
\newcommand{\ra}{\rightarrow}
\newcommand{\phin}{\varphi}
\newcommand{\diver}{\mbox{div}}

\def\eps{{\varepsilon}}
\def\i{{\sl i\;}}
\def\e{{\rm e}}
\def\D{{\rm D\;}}
\def\norm{|\!|\!|}

\DeclareMathOperator{\spann}{span}
\DeclareMathOperator{\cond}{cond}
\DeclareMathOperator{\round}{round}
\DeclareMathOperator{\supp}{supp}

\begin{document}

\parbox{0ex}{ \scalebox{.27}{  \includegraphics[width=.5\textwidth,angle=-90]{cm-logo2.eps} }   }   \\

\parbox{25ex}{
  Prof.~Dr.~S.~Sauter\\
  Institut f"ur Mathematik\\
  Universit"at Z"urich
  }
%
\rule[0cm]{0.cm}{.01cm}                  
\hfill  \parbox{0.6\textwidth}{
  {\sf\LARGE\bfseries Numerik I}
  {\sf\Large\bfseries \;\;--\;\; Homework \blnr }\\[1.5ex]
  Deadline: \sdate\ 13:00
  }



\vspace{10ex}


% \scalebox{.27}{ \put(1250,0){ \includegraphics[width=.5\textwidth,angle=-90]{cm-logo2.eps} }  }


%%%%%%%%%%%%%%%%%%%%%%%%%%%%%%%%%%%%%%%%%%%%%%%%%%%%%%%%%%%%%%%%%%%%%%%%%%%
%
% Aufgaben 
%
%%%%%%%%%%%%%%%%%%%%%%%%%%%%%%%%%%%%%%%%%%%%%%%%%%%%%%%%%%%%%%%%%%%%%%%%%%%

% \vspace{-1cm}
\textbf{Exercise 1} (4 P.)\\
%
Define 
\begin{equation*} y:=333.75b^6+a^2(11a^2b^2-b^6-121b^4-2)+5.5b^8+\frac{a}{2b}.\end{equation*}

\begin {itemize}
\item [(a)] Compute $y$ for $a=\sqrt{4950000001}$ and $b=30000$ using two different tools (e.g. Matlab, Maple/Mathematica). What do you observe? Why?
\item [(b)] The parameters $a=10^6$, $b=10^6+10^{-2}$ are rounded to $a^*$, $b^*$ according to
\begin{align*}
a^*&=a(1+\epsilon_1)\\
b^*&=b(1+\epsilon_2), 
\end{align*}
where $|\epsilon_1|,\, |\epsilon_2| \ll 1$. Use error analysis of first order\footnote{ignoring quadratic terms of \(\epsilon\)} to find $\epsilon_1,\, \epsilon_2$ s.t. $z=a^2-b^2$ and $z^*= a^{*2}-b^{*2}$ satisfy
\begin{equation*}
z^*\doteq z(1+\epsilon_*), 
\quad  |\epsilon_*|\leq 10^{-7}. 
\end{equation*}

\item [(c)] Let 
\begin{equation}\label{1}
y=\cos(x+\delta)-\cos(x),
\end{equation}
 with $\delta>0$ small. Find an equivalent expression of \eqref{1} that avoids round-off errors.
\end{itemize}
 

\textbf{Exercise 2} (5 P.)\\
 We consider the integrals defined by
\begin{equation} I_n = \int_0^1 x^ne^{x} \mathrm{d}x,
\label{intN}
\end{equation}
for  \(n = 0,1,2,3,..\).


\begin{itemize}
\item[(a)] Prove that \begin{equation}\vert I_n\vert\le \dfrac{e}{n+1}.\label{ineq}\end{equation}
\item[(b)] Prove, by using an integration by parts, the recurrence relation
\begin{equation}\begin{array}{c}
I_{k} = e - kI_{k-1}, \quad k \ge 1.
\end{array}\label{Irek}\end{equation}
Compute (with Matlab) $I_{100}$ using formula \eqref{Irek} starting with $I_{0}=e-1.$
Compare your result to the ``exact'' value \(I_{100} = 0.026652359191789498\) and comment.

\item[(c)] Let us introduce $\nu > n, \; \nu \in \mathbb{N}$. Using the reversed (backward) recurrence
\begin{equation}\begin{array}{c}
I_{k} = (e - I_{k+1})/(k+1), \quad k =  \nu - 1, \nu - 2,.., n.
\end{array}\label{IrekRev}\end{equation}
compute (with Matlab) an approximation  $I^*_{100}$ of $I_{100}$ starting with $I^*_{\nu}=0, \text{for } \nu =200.$ Compare the results of (b) and (c).

To explain the differences, derive the relation
\begin{equation}
\delta_n = \left| \dfrac{I_{n} - I_{n}^*}{I_n}\right| \le \left(\prod_{l = n +1}^{\nu} \left|\dfrac{I_{l}}{e - I_{l}}\right|\right) \left|\dfrac{I_{\nu} - I_{\nu}^*}{I_{\nu}} \right|
\end{equation}

and show , using \eqref{ineq}, that  
\begin{equation}
\left(\prod_{l = n+1}^{\nu} \left|\dfrac{I_{l}}{e - I_{l}}\right|\right) \le \left( \dfrac{1}{n}\right)^{\nu - n -1}
\label{relNu}\end{equation}
\item[(d)]  How to choose $\nu$ to satisfy $\vert \delta_{100} \vert\le \epsilon$, for given $\epsilon > 0$?

\end{itemize}


\end{document}